% Part: history
% Chapter: biographies
% Section: julia-robinson

\documentclass[../../../include/open-logic-section]{subfiles}

\begin{document}

\olfileid{his}{bio}{rob}

\olsection{جوليا روبنسون}

\olphoto{robinson-julia}{Julia Robinson}

كانت جوليا بومان روبنسون رياضيةً أمريكية. وهي معروفة أساساً بعملها في
مسائل القرار، وأشهر ما عُرفت به إسهاماتها في حل مسألة هيلبرت العاشرة.
وُلدت روبنسون في سانت لويس بولاية ميزوري في 8 ديسمبر 1919. وتتذكر أنها
كانت مفتونةً بالأعداد منذ طفولتها \citep[4]{Reid1986}. وفي التاسعة من
عمرها أُصيبت بالحمى القرمزية، ثم عانت نوبات متكررة من الحمى الروماتيزمية.
فاضطرت إلى قضاء وقت طويل طريحة الفراش، وتأخرت في تحصيلها الدراسي. ومع أنها
استطاعت تدارك ما فاتها بمساعدة معلمين خصوصيين، تركت الآثار الجسدية لمرضها
أثراً دائماً في حياتها.

على الرغم من صعوبات طفولتها، تخرجت روبنسون في المدرسة الثانوية وقد نالت
جوائز عدة في الرياضيات والعلوم. وبدأت دراستها الجامعية في كلية ولاية سان
دييغو، ثم انتقلت في سنتها الأخيرة إلى جامعة كاليفورنيا في بركلي. وهناك
تأثرت بالرياضي رافائيل روبنسون. وأصبحا صديقين حميمين، وتزوجا عام 1941.
ولكونها زوجةَ عضو في هيئة التدريس، مُنعت روبنسون من التدريس في قسم
الرياضيات في بركلي. وعلى الرغم من استمرارها في حضور مقررات الرياضيات بصفة
مستمعة، كانت تأمل أن تغادر الجامعة وتؤسس أسرة. لكنها أُصيبت بالتهاب رئوي
بعد زواجها بقليل. وأُبلغت بوجود تراكم كبير من النسيج الندبي في قلبها بسبب
الحمى الروماتيزمية التي أصابتها طفلةً. ولشدة ذلك النسيج الندبي، توقع
الطبيب ألا تعيش بعد سن الأربعين، ونصحها بألا تنجب
\citep[13]{Reid1986}.

ظلت روبنسون مكتئبةً مدة طويلة، لكنها قررت في نهاية المطاف مواصلة دراسة
الرياضيات. فعادت إلى بركلي وأتمت الدكتوراه عام 1948 بإشراف ألفرد تارسكي.
كان تارسكي قد أثبت قابلية نظرية الرتبة الأولى للأعداد الحقيقية للقرار،
وكان يلزم من عمل غودل أن نظرية الرتبة الأولى للأعداد الطبيعية غير قابلة
للقرار. وظلت قابلية نظرية الرتبة الأولى للأعداد النسبية للقرار مسألةً
كبرى مفتوحة. وفي أطروحتها \citeyearpar{Robinson1949} برهنت روبنسون أنها
غير قابلة للقرار.

اتجه اهتمام روبنسون إلى مسائل القرار، فحاولت بعد ذلك إيجاد حل لمسألة
هيلبرت العاشرة. وكانت هذه واحدةً من قائمة شهيرة تضم 23 مسألة رياضية طرحها
ديفيد هيلبرت عام 1900. وتسأل المسألة العاشرة عما إذا كانت توجد خوارزمية
تقرر، في زمن منتهٍ، ما إذا كانت لمعادلة كثيرة حدود ذات معاملات صحيحة، مثل
$3x^2 - 2y +3 = 0$، حلول في الأعداد الصحيحة. وتُعرف هذه الأسئلة باسم
\emph{المسائل الديوفانتية}. وبعد نجاحات أولية، انضمت روبنسون إلى مارتن
ديفيس وهيلاري بوتنام، وكانا يعملان أيضاً على المسألة. ونجحوا في بيان أن
المسائل الديوفانتية الأسية (حيث يجوز أن تظهر المجاهيل أسساً أيضاً) غير
قابلة للقرار، وبينوا أن حدساً معيناً (سُمّي لاحقاً «J.R.») يستلزم أن
مسألة هيلبرت العاشرة غير قابلة للقرار \citep{DavisPutnamRobinson1961}.
وواصلت روبنسون العمل على المسألة طوال ستينيات القرن العشرين. وفي عام 1970
برهن الرياضي الروسي الشاب يوري ماتيياسيفيتش أخيراً فرضية J.R. وتسمى
النتيجة المركبة الآن مبرهنة ماتيياسيفيتش--روبنسون--ديفيس--بوتنام، أو
مبرهنة MRDP اختصاراً. وأصبح ماتيياسيفيتش وروبنسون صديقين وتعاونا في أبحاث
عدة. وكتبت روبنسون مرة في رسالة إليه: «يسعدني حقاً أننا، إذ نعمل معاً
(وبيننا آلاف الأميال)، نحرز بوضوح تقدماً أكبر مما كان يمكن لأي منا إحرازه
منفرداً» \citep[45]{Matijasevich1992}.

كانت روبنسون أول امرأة تترأس الجمعية الرياضية الأمريكية، وأول امرأة
تُنتخب في الأكاديمية الوطنية للعلوم (National Academy of Sciences).
وتوفيت في 30 يوليو 1985 عن عمر 65 عاماً بعد تشخيص إصابتها بسرطان الدم.

\begin{reading}
تتوفر أبحاث روبنسون الرياضية في \textit{أعمالها المجموعة}
\citep{Robinson1996}، التي تضم أيضاً إعادة طبع مذكرتها عن السيرة المنشورة
في الأكاديمية الوطنية للعلوم \citep{Feferman1994}. ونشرت شقيقتها الكبرى
كونستانس ريد «سيرة جوليا الذاتية»، استناداً إلى مقابلات
\citep{Reid1986}، فضلاً عن مذكرات كاملة \citep{Reid1996}. وأخرج جورج
تشيشتشيري فيلماً وثائقياً قصيراً عن روبنسون ومسألة هيلبرت العاشرة
\citep{Csicsery2016}. ولمذكرة موجزة عن تعاون يوري ماتيياسيفيتش مع روبنسون
وتأثيرها في عمله، انظر \citep{Matijasevich1992}.
\end{reading}

\end{document}
